\chapter{Unstructured Interviews \label{appendix:c}}

As part of the research methodology, two unstructured interviews were used to triangulated the data generated during the observations period of the case study with the Industry Project
course.

\section{Questionnaire for the Interviews with Students \label{sec:C-1}}

The first interview focused on the reported experiences of the students about the work done with their groups, how did they experience the development of the collaborative work needed for
the course and the mechanisms they perceived were used for team management and the execution of the project. The following was the questionnaire used for these interviews.

\textbf{Question 1.} When working with your current team, either in the classroom or online, where there any \emph{rules} that the group had to follow for any aspect of the work done? Can
you describe them?

\textbf{Question 2.} Can you describe how or when did these rules appear? Had you any influence in the implementation of these rules?

\textbf{Question 3.} Beside your professional role in the course itself (Quantity Surveyor, Construction Manager, Construction Engineer, Architectural Engineer), had you any other role in
your group, related either to the tasks needed for each assignment or for the meetings with your group? Did anyone else in the group had such a role?

\textbf{Question 4.} Please, describe an example of something that went well with the group, something that was helpful in planning or completing a task or assignment.

\textbf{Question 5.} Please, describe an example of something that didn't work with your group, any incident or situation that hindered your group or the expected performance.

\textbf{Question 6.} Did your group manage to solve or circumvent that problem? How?

\textbf{Question 7.} Would you have preferred to complete this project individually? Or did you gain some form of value by working in a team?

This questionnaire was adapted to fit the needs of the project from the original work found in:

\begin{quote}
    Ann Noble in 2000, provided as a supplementary document to The University of Adelaide (2000). Leap into...Collaborative Learning. Centre for Learning and Professional Development.
    Available from \url{http://www.adelaide.edu.au/clpd/resources/leap/}.
\end{quote}

\section{Questionnaire for the Interviews with the Teaching Staff \label{sec:C-2}}

The second interview focused on the views of the teaching staff about how to best promoted good collaborative practices among the students and how this type of collaborative design affects
positively and negatively the execution of the course. The following was the questionnaire used for these interviews.

\textbf{Question 1.} How would you describe a group using good teamwork in the course?

\begin{itemize}
    \item What actions are they taking?
    \item How does it reflect in their work?
\end{itemize}

\textbf{Question 2.} How would you describe a group not succeeding in their teamwork?

\begin{itemize}
    \item Any common bad habits or pitfalls?
    \item How does it reflect in their work?
\end{itemize}

\textbf{Question 3.} Can you identify any mayor difficulties of “pushing” students into proper teamwork?

\textbf{Question 4.} Has collaborative learning help in the overall goals of the ENBU 79X/89X?

\begin{itemize}
    \item Any reasons for the first assignment to be individual?
    \item Good reasons for all assignments to be individual?
    \item Anything you would like to do different in the future?
\end{itemize}

\textbf{Question 5.} How do you deal with students that are reluctant or adverse to group work?

\begin{itemize}
    \item Has the case of Brandon made you change something in the future?
    \item Any other fringe cases you ae aware of? How you delt with it?
\end{itemize}

\textbf{Question 6.} Does the focus on collaborative work eases any aspects of your work as a lecturer?

\textbf{Question 7.} Does the focus on collaborative work makes any aspects of your work as a lecturer harder?

This interview was adapted from questions and methodologies used by \textcite{chiriac_2012,gillies_2010}.